\section{The forensic instrument}
\label{sec:instrument}
\subsection{From requests to architectures: the hypothesis graph}
\label{sec:graph}
The instrument works on one object, the \emph{architecture hypothesis graph} $G_A$.  It
is built by a simple loop that anyone can audit:
\begin{center}
\textbf{request} $\rightarrow$ \textbf{observation} $\rightarrow$ \textbf{constraint}
$\rightarrow$ \textbf{candidate architectures} $\rightarrow$ \textbf{next experiment}.
\end{center}
We send a controlled request.  It might be a pair of prompts that differ by one known
string, so the token-count difference reveals how the tokenizer cuts text.  It might be
a timed decode at a known context length.  It might be a request that deliberately asks
for more output than the cap allows.  Each request returns an observation.  Each
observation is then turned into a \emph{constraint} on the configuration.  When the
observation model justifies it, the constraint is hard: a compatibility restriction
$\mathcal K_i$ that rules candidates out entirely.  An exact count match or a sharp rank
cliff are examples.  Otherwise the constraint is soft: a likelihood factor that ranks
candidates without eliminating any, which is how a timing regime is treated.  Hard
constraints shrink the candidate set, $\mathcal H_t^{\rm hard}=\mathcal
H_0\cap\bigcap_{i\le t}\mathcal K_i$.  Soft evidence ranks whatever survives, $\hat
A_t=\arg\max_{A\in\mathcal H_t^{\rm hard}}S(A\mid D_t)$.  In words: start from every
architecture we know of, remove the ones the hard facts rule out, then score the rest.

The nodes of the graph are the components of the system: tokenizer, core family,
attention or state regime, KV sharing, position scheme, feed-forward topology, output
head, numeric regime, and serving topology.  Each node holds the best-supported
candidate, the full ranked list of candidates, the measured quantities, an evidence
grade, provenance, and the \emph{next experiment} that would separate the survivors.
The edges are functional dependencies, each tagged with the experiment that produced
it.  The loop then spends the next request wherever the remaining ambiguity is largest
(Sec.~\ref{sec:planner}).

This graph is the unit of reconstruction.  It never collapses an endpoint into a single
label, and it never disguises a prior as a measurement.  Where nothing can be measured,
the node says so, lists the candidates, and names the experiment that would decide
between them.

\begin{figure}[t]
\begin{quote}\small
\textbf{Algorithm (the working progression, frozen per run).}
\begin{enumerate}[nosep,leftmargin=1.6em]
\item \textbf{Discover the contract}: send a minimal request; then controlled invalid
      and over-limit variants; read the capabilities and error schema (P01/P05).
\item \textbf{Exploit the cheap channels}: usage accounting (P02 tokenisation deltas),
      wire numerics (P04), determinism (P10).
\item \textbf{Test the generation contract}: continuation length, whether latency grows
      with output length, streaming granularity.  A rejected cap is itself evidence of an
      encoder (P19).
\item \textbf{Probe structure}: output geometry within the access regime (P03);
      decode-state slope (P23); needle retrieval and prefill scaling (P06/P07).
\item \textbf{Follow only unresolved branches}: each U node names its next
      discriminator; the planner prices it by expected information per request.
\item \textbf{Stop at the boundary}: an L5 witness, a $\hat\kappa$-priced cost, or a
      named experiment.  Never a made-up label.
\end{enumerate}
\end{quote}
\caption{The operational progression.  Every step names its failure mode.  Nothing is
asserted without a measurement that produced it.}
\label{fig:algorithm}
\end{figure}

\subsection{Probes and their failure modes}
Table~\ref{tab:probes} organises the probe library by channel.  Each probe declares an
outcome model and returns three things: likelihood evidence, estimates with their
assumptions spelled out, and a named falsifier (the observation that would prove it
wrong).  The full library is in the artifact.  What matters for the claims in this paper
is in the table, including each probe's failure mode as \emph{measured} on \benchname{}.

\begin{table}[t]\centering\small
\caption{Probe families and the architectural constraint each one carries.  The
``remaining alternatives'' travel with the claim into the graph, so a reader always sees
what the probe could not rule out.}
\label{tab:probes}
\begin{tabularx}{\textwidth}{@{}p{1.9cm} Y Y@{}}
\toprule
\textbf{Probe} & \textbf{Constraint produced} & \textbf{Known failure mode (measured)}\\
\midrule
P01/P05 & accepted-input boundary, output cap, error schema & a deployment cap is not the model's capability\\
P02 tokenizer & count-family match set, template overhead & finite strings cannot separate variants that share merge tables (Prop.~\ref{prop:tok})\\
P03 geometry & rank of the common-logit matrix $\Rightarrow$ bottleneck estimate or bound & top-5-only access caps the rank at the number of columns seen (Prop.~\ref{prop:rank})\\
P04/P15 & wire-grid / sampler signatures & rounding in transport is not the compute dtype\\
P06/P07 & retrieval range, prefill/decode scaling & prefill curvature is invisible below about 16k context (Fig.~\ref{fig:attncontrol})\\
P23 & KV-growth vs constant-state slope (streaming per-token times or prefill-cancelling pairs) & drowned in noise on high-overhead production paths (reported U, not wrong)\\
P08--P11 & drafting/cache/jitter/domain-regime signatures & domain timing is the weakest detector (\MoeAcc)\\
P12--P14 & capability/cutoff/format signatures & behaviour is not parameter count\\
P16--P22 & routing/invariance/branch semantics & an unchanged API can hide a model swap (Sec.~\ref{sec:blindfixtures})\\
\bottomrule
\end{tabularx}
\end{table}

\paragraph{The discipline of negative inference.}
The instrument is built around what a channel does \emph{not} say, because that is where
informal probing goes wrong.  Recall is not attention complexity: a wide sliding window
and full attention both retrieve a needle.  Latency is not kernel complexity: queueing
time is not compute time.  Behaviour is not a training recipe: post-training can rewrite
behaviour without touching the architecture.  A candidate that comes only from the prior
is not recovered structure.  And an unchanged API is not an unchanged model.  Every
evidence record carries its alternatives, and every figure draws prior-only values in
grey.

\subsection{Evidence combination that respects dependence}
\label{sec:grouped}
Each piece of evidence is a likelihood factor.  For property $j$ and candidate value
$v$, the support is
$q_j(v\mid D)\propto q_{j,0}(v)\prod_e\max(\Lambda_e(v),\epsilon)^{s_e}$.
Read it this way: start from the prior $q_{j,0}(v)$, then for each evidence unit $e$
multiply by that unit's likelihood $\Lambda_e(v)$, raised to a strength $s_e$ that says
how much we trust the unit.  The floor $\epsilon$ stops any single factor from driving a
candidate to exactly zero.  Version 2 of this paper multiplied every factor
independently.  That has a hidden flaw: a long series of similar probes gets counted
again and again, and confidence inflates by repetition.  The benchmark let us measure
the damage.  We now combine evidence \emph{group-aware}: factors that come from the same
experiment group (one probe series, one randomised block) are averaged into a single
update.  We report calibration for both schemes on ground truth
(Table~\ref{tab:calibration}).  Scores are still \emph{support}, not probabilities, but
the benchmark now tells us how far from a probability they are (expected calibration
error \BenchECE{}) and where.

Two views accompany every posterior in the paper.  The \emph{reference-prior} view uses
priors that encode the 2024--2026 model landscape.  The \emph{evidence-only} view starts
from a uniform prior.  Showing both separates what the data did from what the prior
did.  The gap matters most exactly where a strong prior meets a single decisive channel
(Sec.~\ref{sec:familysplit}).

\subsection{The planner}
\label{sec:planner}
The planner chooses the next probe by expected information gain per predicted request.
In plain terms: for each probe it could run, it estimates how much the uncertainty
would shrink (weighted by how diagnostic the probe is), divides by how many requests
the probe would cost, and picks the best ratio.  It is a heuristic scheduler, not an
optimal Bayesian design.  Sec.~\ref{sec:plannerlive} evaluates it \emph{live} against
the simulator (not from replayed logs) with request budgets.  That study found a real
pathology, front-runner starvation, and the fix we adopted (the core-coverage floor).
